\section*{Reproducibility Statement}
\label{sec:reproducibility}

We detail the steps taken to ensure full reproducibility of our results.

\paragraph{Code and Data.}
All source code, synthetic data generators, experiment scripts, and analysis pipelines are publicly available at \url{https://github.com/VJlaxmi/movie-scheduling-paper}.
The repository includes:
\begin{itemize}
    \item Complete BDM-AOS implementation with all ablation variants ({\ttfamily bdm\_aos/} package);
    \item Baseline implementations (GA, PSO, MILP) using the same evaluation framework;
    \item Synthetic instance generator with fixed seed ({\ttfamily seed=42}) for deterministic instance creation;
    \item Experiment runner with 30 independent seeds per algorithm-instance pair;
    \item Analysis scripts that reproduce all tables and figures from raw JSON results.
\end{itemize}

\paragraph{Instance Generation.}
All 8 benchmark instances (S10--S100) are generated deterministically from a single master seed.
Instance parameters---including the three-tier location hierarchy (cities, towns, locations), actor criticality weights, location criticality weights, maximum consecutive working days, and hierarchical transfer costs---are documented in Table~\ref{tab:instances} and Appendix~\ref{app:instances}.
The generator produces consistent actor availability calendars, location compatibility matrices, scene-actor assignment matrices, town-location mappings, and all parameters required for the 13 constraint families of Section~\ref{sec:formulation}.

\paragraph{Algorithm Parameters.}
All algorithm hyperparameters are fixed throughout the experiments and documented in Section~\ref{sec:experiments}.
BDM-AOS uses: population size $N=20$, generations $G=40$, crossover rate $p_c=0.9$, mutation rate $p_m=0.3$, MILP sub-problem time limit 10s.
GA uses: population size 50, 100 generations, $p_c=0.85$, $p_m=0.15$.
PSO uses: swarm size 30, 100 iterations.
MILP uses: 60-second time limit with CBC solver.

\paragraph{Statistical Protocol.}
Each algorithm is run 30 times with seeds 1--30 on each instance.
We report mean $\pm$ standard deviation for cost and makespan.
Statistical significance is assessed using Friedman tests with Holm's post-hoc correction at significance level $\alpha = 0.05$.
Quality indicators (hypervolume, inverted generational distance) use a common reference point derived from the worst values across all algorithms.

\paragraph{Computational Environment.}
All experiments were run on a single machine with Python 3.14 and the following key libraries: NumPy, SciPy, pandas, NetworkX, python-louvain (community detection), Pyomo with CBC solver, and Matplotlib.
Runtime measurements use Python's {\ttfamily time.time()} for wall-clock timing.
